%%%%%%%%%%%%%%%%%%%%%%%%%%%%%%%%%%%%%%%%%%%%%%%%%%%%%%%%%%%%
% 10.13 INVENTORY THEORY
% The Composition of Advanced Mathematics
%%%%%%%%%%%%%%%%%%%%%%%%%%%%%%%%%%%%%%%%%%%%%%%%%%%%%%%%%%%%

\section{Inventory Theory}
\label{sec:inventory-theory}

\prerequisite{
Optimization fundamentals, calculus, probability, random variables,
expectation, stochastic optimization, dynamic programming, queueing
theory, and basic operations research.
}

\learningobjectives{
By the end of this section, the reader should be able to:
\begin{itemize}
    \item define inventory and identify its operational purposes;
    \item distinguish cycle stock, safety stock, pipeline inventory,
          anticipation inventory, and decoupling inventory;
    \item identify ordering, holding, shortage, and purchase costs;
    \item distinguish deterministic and stochastic inventory models;
    \item distinguish continuous-review and periodic-review systems;
    \item formulate the classical Economic Order Quantity model;
    \item derive the EOQ formula;
    \item compute optimal cycle inventory and order frequency;
    \item understand the cost behavior around EOQ;
    \item formulate the Economic Production Quantity model;
    \item analyze finite replenishment rates;
    \item model quantity discounts;
    \item compute reorder points;
    \item incorporate deterministic lead time;
    \item define lead-time demand;
    \item understand safety stock;
    \item distinguish cycle-service level and fill rate;
    \item compute safety stock under normal demand assumptions;
    \item formulate continuous-review \((Q,r)\) policies;
    \item formulate periodic-review order-up-to policies;
    \item understand the newsvendor problem;
    \item derive the critical fractile;
    \item distinguish underage and overage costs;
    \item solve discrete and continuous newsvendor models;
    \item understand lost-sales and backorder systems;
    \item recognize base-stock policies;
    \item understand inventory-position concepts;
    \item formulate finite-horizon inventory dynamic programs;
    \item understand the structure of \((s,S)\) policies conceptually;
    \item recognize multi-item inventory constraints;
    \item understand ABC inventory classification conceptually;
    \item recognize multi-echelon inventory systems;
    \item understand the bullwhip effect conceptually;
    \item connect inventory theory with supply chains, manufacturing,
          retail, logistics, healthcare, maintenance, and finance.
\end{itemize}
}

%%%%%%%%%%%%%%%%%%%%%%%%%%%%%%%%%%%%%%%%%%%%%%%%%%%%%%%%%%%%
\subsection{What Is Inventory?}
\label{subsec:what-is-inventory}
%%%%%%%%%%%%%%%%%%%%%%%%%%%%%%%%%%%%%%%%%%%%%%%%%%%%%%%%%%%%

Inventory is material, product, or capacity held for future use.

Examples include:

\[
\boxed{
\text{raw materials},
}
\]

\[
\boxed{
\text{work in process},
}
\]

\[
\boxed{
\text{finished goods},
}
\]

\[
\boxed{
\text{spare parts},
}
\]

and

\[
\boxed{
\text{supplies}.
}
\]

Inventory allows supply and demand to occur at different times.

%%%%%%%%%%%%%%%%%%%%%%%%%%%%%%%%%%%%%%%%%%%%%%%%%%%%%%%%%%%%
\subsection{Why Hold Inventory?}
\label{subsec:why-hold-inventory}
%%%%%%%%%%%%%%%%%%%%%%%%%%%%%%%%%%%%%%%%%%%%%%%%%%%%%%%%%%%%

Inventory can be used to:

\[
\boxed{
\text{buffer demand uncertainty},
}
\]

\[
\boxed{
\text{buffer supply uncertainty},
}
\]

\[
\boxed{
\text{reduce ordering frequency},
}
\]

\[
\boxed{
\text{support production continuity},
}
\]

and

\[
\boxed{
\text{prepare for expected future demand}.
}
\]

However, inventory also creates cost.

%%%%%%%%%%%%%%%%%%%%%%%%%%%%%%%%%%%%%%%%%%%%%%%%%%%%%%%%%%%%
\subsection{Types of Inventory}
\label{subsec:types-of-inventory}
%%%%%%%%%%%%%%%%%%%%%%%%%%%%%%%%%%%%%%%%%%%%%%%%%%%%%%%%%%%%

Common inventory categories include:

\[
\boxed{
\text{cycle stock},
}
\]

\[
\boxed{
\text{safety stock},
}
\]

\[
\boxed{
\text{pipeline inventory},
}
\]

\[
\boxed{
\text{anticipation inventory},
}
\]

and

\[
\boxed{
\text{decoupling inventory}.
}
\]

%%%%%%%%%%%%%%%%%%%%%%%%%%%%%%%%%%%%%%%%%%%%%%%%%%%%%%%%%%%%
\subsection{Cycle Stock}
\label{subsec:cycle-stock}
%%%%%%%%%%%%%%%%%%%%%%%%%%%%%%%%%%%%%%%%%%%%%%%%%%%%%%%%%%%%

Cycle stock is inventory created by ordering or producing in batches.

If an order of size \(Q\) arrives instantaneously and demand is constant,
inventory may decline approximately from

\[
Q
\]

to

\[
0.
\]

The average cycle stock is then

\[
\boxed{
\frac{Q}{2}.
}
\]

%%%%%%%%%%%%%%%%%%%%%%%%%%%%%%%%%%%%%%%%%%%%%%%%%%%%%%%%%%%%
\subsection{Safety Stock}
\label{subsec:safety-stock}
%%%%%%%%%%%%%%%%%%%%%%%%%%%%%%%%%%%%%%%%%%%%%%%%%%%%%%%%%%%%

Safety stock is extra inventory held to protect against uncertainty.

It may buffer uncertainty in:

\[
\boxed{
\text{demand},
}
\]

\[
\boxed{
\text{lead time},
}
\]

or both.

Safety stock trades increased holding cost for improved service
reliability.

%%%%%%%%%%%%%%%%%%%%%%%%%%%%%%%%%%%%%%%%%%%%%%%%%%%%%%%%%%%%
\subsection{Pipeline Inventory}
\label{subsec:pipeline-inventory}
%%%%%%%%%%%%%%%%%%%%%%%%%%%%%%%%%%%%%%%%%%%%%%%%%%%%%%%%%%%%

Pipeline inventory is material currently moving through the supply or
production system.

If demand rate is \(D\) units per unit time and replenishment lead time is
\(L\), average demand occurring during lead time is

\[
\boxed{
DL.
}
\]

This quantity is closely connected with the reorder point.

%%%%%%%%%%%%%%%%%%%%%%%%%%%%%%%%%%%%%%%%%%%%%%%%%%%%%%%%%%%%
\subsection{Inventory Costs}
\label{subsec:inventory-costs}
%%%%%%%%%%%%%%%%%%%%%%%%%%%%%%%%%%%%%%%%%%%%%%%%%%%%%%%%%%%%

A basic inventory model may include:

\[
\boxed{
\text{purchase cost},
}
\]

\[
\boxed{
\text{ordering or setup cost},
}
\]

\[
\boxed{
\text{holding cost},
}
\]

and

\[
\boxed{
\text{shortage cost}.
}
\]

The optimal inventory policy balances these competing costs.

%%%%%%%%%%%%%%%%%%%%%%%%%%%%%%%%%%%%%%%%%%%%%%%%%%%%%%%%%%%%
\subsection{Ordering Cost}
\label{subsec:ordering-cost}
%%%%%%%%%%%%%%%%%%%%%%%%%%%%%%%%%%%%%%%%%%%%%%%%%%%%%%%%%%%%

Let

\[
\boxed{
K
=
\text{fixed cost per replenishment order}.
}
\]

If annual demand is \(D\) and each order contains \(Q\) units, the number
of orders per year is

\[
\boxed{
\frac{D}{Q}.
}
\]

Therefore annual ordering cost is

\[
\boxed{
K\frac{D}{Q}.
}
\]

%%%%%%%%%%%%%%%%%%%%%%%%%%%%%%%%%%%%%%%%%%%%%%%%%%%%%%%%%%%%
\subsection{Holding Cost}
\label{subsec:holding-cost}
%%%%%%%%%%%%%%%%%%%%%%%%%%%%%%%%%%%%%%%%%%%%%%%%%%%%%%%%%%%%

Let

\[
\boxed{
h
=
\text{holding cost per unit per unit time}.
}
\]

If average inventory is

\[
\frac{Q}{2},
\]

annual holding cost is

\[
\boxed{
h\frac{Q}{2}.
}
\]

Holding cost may reflect:

\[
\boxed{
\text{capital cost},
}
\]

\[
\boxed{
\text{storage},
}
\]

\[
\boxed{
\text{insurance},
}
\]

\[
\boxed{
\text{obsolescence},
}
\]

and

\[
\boxed{
\text{damage or shrinkage}.
}
\]

%%%%%%%%%%%%%%%%%%%%%%%%%%%%%%%%%%%%%%%%%%%%%%%%%%%%%%%%%%%%
\subsection{Shortage Cost}
\label{subsec:shortage-cost}
%%%%%%%%%%%%%%%%%%%%%%%%%%%%%%%%%%%%%%%%%%%%%%%%%%%%%%%%%%%%

A shortage occurs when demand cannot be immediately satisfied from stock.

Possible consequences include:

\[
\boxed{
\text{backorder cost},
}
\]

\[
\boxed{
\text{lost sales},
}
\]

\[
\boxed{
\text{expediting cost},
}
\]

or

\[
\boxed{
\text{service penalties}.
}
\]

Shortage modeling depends on whether unmet demand is delayed or lost.

%%%%%%%%%%%%%%%%%%%%%%%%%%%%%%%%%%%%%%%%%%%%%%%%%%%%%%%%%%%%
\subsection{Deterministic and Stochastic Inventory}
\label{subsec:deterministic-stochastic-inventory}
%%%%%%%%%%%%%%%%%%%%%%%%%%%%%%%%%%%%%%%%%%%%%%%%%%%%%%%%%%%%

A deterministic model assumes quantities such as demand and lead time are
known.

A stochastic model treats one or more quantities as random.

Thus:

\[
\boxed{
\text{deterministic inventory}
\rightarrow
\text{known demand},
}
\]

while

\[
\boxed{
\text{stochastic inventory}
\rightarrow
\text{uncertain demand or supply}.
}
\]

%%%%%%%%%%%%%%%%%%%%%%%%%%%%%%%%%%%%%%%%%%%%%%%%%%%%%%%%%%%%
\subsection{Continuous and Periodic Review}
\label{subsec:continuous-periodic-review}
%%%%%%%%%%%%%%%%%%%%%%%%%%%%%%%%%%%%%%%%%%%%%%%%%%%%%%%%%%%%

In continuous review, inventory status is monitored continuously.

A replenishment may be triggered whenever inventory position reaches a
specified level.

In periodic review, inventory is checked only at scheduled review times:

\[
\boxed{
0,R,2R,3R,\ldots.
}
\]

The two systems require different protection against uncertainty.

%%%%%%%%%%%%%%%%%%%%%%%%%%%%%%%%%%%%%%%%%%%%%%%%%%%%%%%%%%%%
\subsection{Economic Order Quantity Model}
\label{subsec:economic-order-quantity}
%%%%%%%%%%%%%%%%%%%%%%%%%%%%%%%%%%%%%%%%%%%%%%%%%%%%%%%%%%%%

The classical Economic Order Quantity model, abbreviated EOQ, balances
ordering cost and holding cost.

Its standard assumptions include:

\begin{itemize}
    \item constant deterministic demand rate;
    \item instantaneous replenishment;
    \item no shortages;
    \item constant ordering cost;
    \item constant holding cost;
    \item one item;
    \item no relevant quantity discount in the basic model.
\end{itemize}

%%%%%%%%%%%%%%%%%%%%%%%%%%%%%%%%%%%%%%%%%%%%%%%%%%%%%%%%%%%%
\subsection{EOQ Cost Function}
\label{subsec:eoq-cost-function}
%%%%%%%%%%%%%%%%%%%%%%%%%%%%%%%%%%%%%%%%%%%%%%%%%%%%%%%%%%%%

Let:

\[
D
=
\text{demand per unit time},
\]

\[
K
=
\text{ordering cost per order},
\]

\[
h
=
\text{holding cost per unit per unit time}.
\]

Ignoring purchase cost that is independent of \(Q\), the relevant cost is

\[
\boxed{
TC(Q)
=
K\frac{D}{Q}
+
h\frac{Q}{2}.
}
\]

%%%%%%%%%%%%%%%%%%%%%%%%%%%%%%%%%%%%%%%%%%%%%%%%%%%%%%%%%%%%
\subsection{Derivation of EOQ}
\label{subsec:eoq-derivation}
%%%%%%%%%%%%%%%%%%%%%%%%%%%%%%%%%%%%%%%%%%%%%%%%%%%%%%%%%%%%

Differentiate:

\[
\frac{d}{dQ}
TC(Q)
=
-\frac{KD}{Q^2}
+
\frac{h}{2}.
\]

Set the derivative equal to zero:

\[
-\frac{KD}{Q^2}
+
\frac{h}{2}
=
0.
\]

Thus,

\[
\frac{KD}{Q^2}
=
\frac{h}{2}.
\]

Therefore,

\[
Q^2
=
\frac{2KD}{h}.
\]

Hence,

\[
\boxed{
Q^*
=
\sqrt{
\frac{2KD}{h}
}.
}
\]

%%%%%%%%%%%%%%%%%%%%%%%%%%%%%%%%%%%%%%%%%%%%%%%%%%%%%%%%%%%%
\subsection{Second-Derivative Verification}
\label{subsec:eoq-second-derivative}
%%%%%%%%%%%%%%%%%%%%%%%%%%%%%%%%%%%%%%%%%%%%%%%%%%%%%%%%%%%%

The second derivative is

\[
\frac{d^2}{dQ^2}
TC(Q)
=
\frac{2KD}{Q^3}.
\]

For

\[
Q>0,
\]

\[
\boxed{
\frac{d^2}{dQ^2}TC(Q)>0.
}
\]

Therefore the EOQ stationary point is a strict minimum.

%%%%%%%%%%%%%%%%%%%%%%%%%%%%%%%%%%%%%%%%%%%%%%%%%%%%%%%%%%%%
\subsection{Ordering and Holding Cost at EOQ}
\label{subsec:eoq-cost-balance}
%%%%%%%%%%%%%%%%%%%%%%%%%%%%%%%%%%%%%%%%%%%%%%%%%%%%%%%%%%%%

At the optimum,

\[
\frac{KD}{(Q^*)^2}
=
\frac{h}{2}.
\]

Multiplying by \(Q^*\),

\[
\boxed{
K\frac{D}{Q^*}
=
h\frac{Q^*}{2}.
}
\]

Thus annual ordering cost equals annual cycle-stock holding cost at EOQ.

%%%%%%%%%%%%%%%%%%%%%%%%%%%%%%%%%%%%%%%%%%%%%%%%%%%%%%%%%%%%
\subsection{Minimum Relevant EOQ Cost}
\label{subsec:eoq-minimum-cost}
%%%%%%%%%%%%%%%%%%%%%%%%%%%%%%%%%%%%%%%%%%%%%%%%%%%%%%%%%%%%

At

\[
Q^*
=
\sqrt{\frac{2KD}{h}},
\]

the minimum relevant cost is

\[
\boxed{
TC(Q^*)
=
\sqrt{2KDh}.
}
\]

This excludes any purchase cost independent of order size.

%%%%%%%%%%%%%%%%%%%%%%%%%%%%%%%%%%%%%%%%%%%%%%%%%%%%%%%%%%%%
\subsection{Worked Example: EOQ}
\label{subsec:worked-eoq}
%%%%%%%%%%%%%%%%%%%%%%%%%%%%%%%%%%%%%%%%%%%%%%%%%%%%%%%%%%%%

\begin{workedexamplebox}{Optimal Order Quantity}

Suppose annual demand is

\[
D=2400
\]

units.

Ordering cost is

\[
K=50
\]

per order.

Annual holding cost is

\[
h=6
\]

per unit.

Then

\[
Q^*
=
\sqrt{
\frac{
2(50)(2400)
}{
6
}
}.
\]

Thus,

\[
Q^*
=
\sqrt{40000}
=
200.
\]

\begin{answerbox}
\[
\boxed{
Q^*=200\text{ units}.
}
\]
\end{answerbox}

The number of orders per year is

\[
\frac{2400}{200}
=
12.
\]

Thus the firm orders approximately once per month.

\end{workedexamplebox}

%%%%%%%%%%%%%%%%%%%%%%%%%%%%%%%%%%%%%%%%%%%%%%%%%%%%%%%%%%%%
\subsection{EOQ Cycle Length}
\label{subsec:eoq-cycle-length}
%%%%%%%%%%%%%%%%%%%%%%%%%%%%%%%%%%%%%%%%%%%%%%%%%%%%%%%%%%%%

If demand rate is \(D\), an order of size \(Q\) lasts

\[
\boxed{
T
=
\frac{Q}{D}.
}
\]

At EOQ,

\[
\boxed{
T^*
=
\frac{Q^*}{D}.
}
\]

The order frequency is

\[
\boxed{
\frac{1}{T^*}
=
\frac{D}{Q^*}.
}
\]

%%%%%%%%%%%%%%%%%%%%%%%%%%%%%%%%%%%%%%%%%%%%%%%%%%%%%%%%%%%%
\subsection{EOQ Cost Robustness}
\label{subsec:eoq-cost-robustness}
%%%%%%%%%%%%%%%%%%%%%%%%%%%%%%%%%%%%%%%%%%%%%%%%%%%%%%%%%%%%

The EOQ cost function is relatively flat near its minimum.

Therefore modest errors in \(Q\) may produce only small cost increases.

This means practical order quantities can often be rounded to convenient
values without a large cost penalty.

%%%%%%%%%%%%%%%%%%%%%%%%%%%%%%%%%%%%%%%%%%%%%%%%%%%%%%%%%%%%
\subsection{Economic Production Quantity}
\label{subsec:economic-production-quantity}
%%%%%%%%%%%%%%%%%%%%%%%%%%%%%%%%%%%%%%%%%%%%%%%%%%%%%%%%%%%%

The Economic Production Quantity model, abbreviated EPQ, assumes
replenishment occurs gradually rather than instantaneously.

Let:

\[
p
=
\text{production rate},
\]

and

\[
d
=
\text{demand rate},
\]

with

\[
\boxed{
p>d.
}
\]

During production, inventory increases at net rate

\[
\boxed{
p-d.
}
\]

%%%%%%%%%%%%%%%%%%%%%%%%%%%%%%%%%%%%%%%%%%%%%%%%%%%%%%%%%%%%
\subsection{Maximum Inventory in EPQ}
\label{subsec:epq-maximum-inventory}
%%%%%%%%%%%%%%%%%%%%%%%%%%%%%%%%%%%%%%%%%%%%%%%%%%%%%%%%%%%%

Producing batch \(Q\) requires time

\[
\boxed{
\frac{Q}{p}.
}
\]

During that time, net inventory accumulates at rate \(p-d\).

Therefore maximum inventory is

\[
\boxed{
I_{\max}
=
\left(
p-d
\right)
\frac{Q}{p}.
}
\]

Thus,

\[
\boxed{
I_{\max}
=
Q
\left(
1-\frac{d}{p}
\right).
}
\]

%%%%%%%%%%%%%%%%%%%%%%%%%%%%%%%%%%%%%%%%%%%%%%%%%%%%%%%%%%%%
\subsection{Average Inventory in EPQ}
\label{subsec:epq-average-inventory}
%%%%%%%%%%%%%%%%%%%%%%%%%%%%%%%%%%%%%%%%%%%%%%%%%%%%%%%%%%%%

Under the standard sawtooth pattern,

\[
\boxed{
I_{\mathrm{avg}}
=
\frac{
I_{\max}
}{
2
}.
}
\]

Hence,

\[
\boxed{
I_{\mathrm{avg}}
=
\frac{Q}{2}
\left(
1-\frac{d}{p}
\right).
}
\]

%%%%%%%%%%%%%%%%%%%%%%%%%%%%%%%%%%%%%%%%%%%%%%%%%%%%%%%%%%%%
\subsection{EPQ Cost Function}
\label{subsec:epq-cost-function}
%%%%%%%%%%%%%%%%%%%%%%%%%%%%%%%%%%%%%%%%%%%%%%%%%%%%%%%%%%%%

The relevant cost becomes

\[
\boxed{
TC(Q)
=
K\frac{D}{Q}
+
h
\frac{Q}{2}
\left(
1-\frac{d}{p}
\right).
}
\]

Differentiating produces the optimal production quantity

\[
\boxed{
Q^*
=
\sqrt{
\frac{
2KD
}{
h
\left(
1-\frac{d}{p}
\right)
}
}.
}
\]

%%%%%%%%%%%%%%%%%%%%%%%%%%%%%%%%%%%%%%%%%%%%%%%%%%%%%%%%%%%%
\subsection{EOQ as a Limit of EPQ}
\label{subsec:eoq-limit-epq}
%%%%%%%%%%%%%%%%%%%%%%%%%%%%%%%%%%%%%%%%%%%%%%%%%%%%%%%%%%%%

If production becomes arbitrarily fast,

\[
p\to\infty,
\]

then

\[
\frac{d}{p}\to0.
\]

Therefore,

\[
Q^*
\to
\sqrt{
\frac{2KD}{h}
}.
\]

Thus EOQ is the instantaneous-replenishment limit of EPQ.

%%%%%%%%%%%%%%%%%%%%%%%%%%%%%%%%%%%%%%%%%%%%%%%%%%%%%%%%%%%%
\subsection{Worked Example: EPQ}
\label{subsec:worked-epq}
%%%%%%%%%%%%%%%%%%%%%%%%%%%%%%%%%%%%%%%%%%%%%%%%%%%%%%%%%%%%

\begin{workedexamplebox}{Finite Production Rate}

Suppose:

\[
D=10000,
\]

\[
K=100,
\]

\[
h=5,
\]

\[
p=50000,
\]

and

\[
d=10000.
\]

Then

\[
1-\frac{d}{p}
=
1-\frac{10000}{50000}
=
0.8.
\]

Therefore,

\[
Q^*
=
\sqrt{
\frac{
2(100)(10000)
}{
5(0.8)
}
}.
\]

Thus,

\[
Q^*
=
\sqrt{500000}
\approx707.1.
\]

\begin{answerbox}
\[
\boxed{
Q^*\approx707\text{ units}.
}
\]
\end{answerbox}

\end{workedexamplebox}

%%%%%%%%%%%%%%%%%%%%%%%%%%%%%%%%%%%%%%%%%%%%%%%%%%%%%%%%%%%%
\subsection{Quantity Discounts}
\label{subsec:quantity-discounts}
%%%%%%%%%%%%%%%%%%%%%%%%%%%%%%%%%%%%%%%%%%%%%%%%%%%%%%%%%%%%

A supplier may reduce unit purchase price when order quantity exceeds
specified thresholds.

Then purchase cost depends on \(Q\).

Total cost becomes

\[
\boxed{
TC(Q)
=
c(Q)D
+
K\frac{D}{Q}
+
h(Q)\frac{Q}{2}.
}
\]

The optimal order quantity may lie at an EOQ candidate or at a discount
breakpoint.

%%%%%%%%%%%%%%%%%%%%%%%%%%%%%%%%%%%%%%%%%%%%%%%%%%%%%%%%%%%%
\subsection{All-Units Quantity Discount}
\label{subsec:all-units-discount}
%%%%%%%%%%%%%%%%%%%%%%%%%%%%%%%%%%%%%%%%%%%%%%%%%%%%%%%%%%%%

In an all-units discount, once a quantity threshold is reached, the lower
price applies to every unit in the order.

Suppose:

\[
c_1
\quad
\text{for}
\quad
Q<q_1,
\]

and

\[
c_2<c_1
\quad
\text{for}
\quad
Q\geq q_1.
\]

The purchase-cost reduction can justify ordering more than the ordinary
EOQ.

%%%%%%%%%%%%%%%%%%%%%%%%%%%%%%%%%%%%%%%%%%%%%%%%%%%%%%%%%%%%
\subsection{Quantity-Discount Procedure}
\label{subsec:quantity-discount-procedure}
%%%%%%%%%%%%%%%%%%%%%%%%%%%%%%%%%%%%%%%%%%%%%%%%%%%%%%%%%%%%

A standard procedure is:

\begin{enumerate}
    \item compute an EOQ for each price level;
    \item check whether each EOQ is feasible in that price interval;
    \item if not, use the nearest feasible breakpoint candidate;
    \item compute total annual cost for every candidate;
    \item choose the lowest-cost feasible candidate.
\end{enumerate}

%%%%%%%%%%%%%%%%%%%%%%%%%%%%%%%%%%%%%%%%%%%%%%%%%%%%%%%%%%%%
\subsection{Reorder Point}
\label{subsec:reorder-point}
%%%%%%%%%%%%%%%%%%%%%%%%%%%%%%%%%%%%%%%%%%%%%%%%%%%%%%%%%%%%

Order quantity answers:

\[
\boxed{
\text{how much to order?}
}
\]

The reorder point answers:

\[
\boxed{
\text{when to order?}
}
\]

With deterministic demand rate \(d\) and lead time \(L\),

\[
\boxed{
r
=
dL.
}
\]

%%%%%%%%%%%%%%%%%%%%%%%%%%%%%%%%%%%%%%%%%%%%%%%%%%%%%%%%%%%%
\subsection{Lead Time}
\label{subsec:inventory-lead-time}
%%%%%%%%%%%%%%%%%%%%%%%%%%%%%%%%%%%%%%%%%%%%%%%%%%%%%%%%%%%%

Lead time is the time between placing an order and receiving the
replenishment.

If demand continues during lead time, enough inventory must remain to
cover that demand.

Hence reorder decisions depend on

\[
\boxed{
\text{lead-time demand}.
}
\]

%%%%%%%%%%%%%%%%%%%%%%%%%%%%%%%%%%%%%%%%%%%%%%%%%%%%%%%%%%%%
\subsection{Worked Example: Deterministic Reorder Point}
\label{subsec:worked-deterministic-reorder-point}
%%%%%%%%%%%%%%%%%%%%%%%%%%%%%%%%%%%%%%%%%%%%%%%%%%%%%%%%%%%%

\begin{workedexamplebox}{Reordering Before Stockout}

Suppose demand is

\[
20
\]

units per day and lead time is

\[
6
\]

days.

Then

\[
r
=
20(6)
=
120.
\]

\begin{answerbox}
\[
\boxed{
r=120\text{ units}.
}
\]
\end{answerbox}

When inventory position reaches \(120\), a replenishment should be
triggered under the deterministic model.

\end{workedexamplebox}

%%%%%%%%%%%%%%%%%%%%%%%%%%%%%%%%%%%%%%%%%%%%%%%%%%%%%%%%%%%%
\subsection{Inventory Position}
\label{subsec:inventory-position}
%%%%%%%%%%%%%%%%%%%%%%%%%%%%%%%%%%%%%%%%%%%%%%%%%%%%%%%%%%%%

Inventory position is commonly defined as

\[
\boxed{
IP
=
\text{on-hand inventory}
+
\text{on-order inventory}
-
\text{backorders}.
}
\]

Continuous-review reorder decisions are typically based on inventory
position rather than on-hand inventory alone.

%%%%%%%%%%%%%%%%%%%%%%%%%%%%%%%%%%%%%%%%%%%%%%%%%%%%%%%%%%%%
\subsection{Why Inventory Position Matters}
\label{subsec:why-inventory-position-matters}
%%%%%%%%%%%%%%%%%%%%%%%%%%%%%%%%%%%%%%%%%%%%%%%%%%%%%%%%%%%%

Suppose on-hand inventory is low but a large replenishment is already on
the way.

Looking only at on-hand inventory could trigger unnecessary duplicate
orders.

Inventory position includes outstanding replenishments and gives a more
complete picture of future stock.

%%%%%%%%%%%%%%%%%%%%%%%%%%%%%%%%%%%%%%%%%%%%%%%%%%%%%%%%%%%%
\subsection{Uncertain Lead-Time Demand}
\label{subsec:uncertain-lead-time-demand}
%%%%%%%%%%%%%%%%%%%%%%%%%%%%%%%%%%%%%%%%%%%%%%%%%%%%%%%%%%%%

If demand is random, lead-time demand is a random variable

\[
\boxed{
D_L.
}
\]

A basic reorder point may be written

\[
\boxed{
r
=
\mathbb{E}[D_L]
+
SS,
}
\]

where

\[
SS
=
\text{safety stock}.
\]

%%%%%%%%%%%%%%%%%%%%%%%%%%%%%%%%%%%%%%%%%%%%%%%%%%%%%%%%%%%%
\subsection{Normally Distributed Lead-Time Demand}
\label{subsec:normal-lead-time-demand}
%%%%%%%%%%%%%%%%%%%%%%%%%%%%%%%%%%%%%%%%%%%%%%%%%%%%%%%%%%%%

Suppose

\[
D_L
\sim
N
(
\mu_L,
\sigma_L^2
).
\]

To achieve a chosen cycle-service probability \(\alpha\), select

\[
\boxed{
r
=
\mu_L
+
z_\alpha\sigma_L,
}
\]

where

\[
z_\alpha
\]

is the standard-normal \(\alpha\)-quantile.

Therefore,

\[
\boxed{
SS
=
z_\alpha\sigma_L.
}
\]

%%%%%%%%%%%%%%%%%%%%%%%%%%%%%%%%%%%%%%%%%%%%%%%%%%%%%%%%%%%%
\subsection{Demand Variability During Constant Lead Time}
\label{subsec:demand-variability-lead-time}
%%%%%%%%%%%%%%%%%%%%%%%%%%%%%%%%%%%%%%%%%%%%%%%%%%%%%%%%%%%%

Suppose daily demands are independent with mean

\[
\mu_D
\]

and standard deviation

\[
\sigma_D,
\]

and lead time is \(L\) days.

Then

\[
\boxed{
\mu_L
=
L\mu_D,
}
\]

and

\[
\boxed{
\sigma_L
=
\sqrt{L}\sigma_D.
}
\]

Therefore,

\[
\boxed{
r
=
L\mu_D
+
z_\alpha
\sqrt{L}\sigma_D.
}
\]

%%%%%%%%%%%%%%%%%%%%%%%%%%%%%%%%%%%%%%%%%%%%%%%%%%%%%%%%%%%%
\subsection{Worked Example: Safety Stock}
\label{subsec:worked-safety-stock}
%%%%%%%%%%%%%%%%%%%%%%%%%%%%%%%%%%%%%%%%%%%%%%%%%%%%%%%%%%%%

\begin{workedexamplebox}{Normally Distributed Lead-Time Demand}

Suppose daily demand has

\[
\mu_D=50,
\]

\[
\sigma_D=8,
\]

and lead time is

\[
L=4
\]

days.

Then

\[
\mu_L
=
4(50)
=
200.
\]

The lead-time standard deviation is

\[
\sigma_L
=
\sqrt{4}(8)
=
16.
\]

For

\[
z=1.645,
\]

corresponding approximately to a \(95\%\) one-sided normal quantile,

\[
SS
=
1.645(16)
=
26.32.
\]

Thus,

\[
r
=
200+26.32
=
226.32.
\]

\begin{answerbox}
\[
\boxed{
SS\approx26.3,
\qquad
r\approx226.3.
}
\]
\end{answerbox}

A practical integer reorder point might be rounded upward as appropriate.

\end{workedexamplebox}

%%%%%%%%%%%%%%%%%%%%%%%%%%%%%%%%%%%%%%%%%%%%%%%%%%%%%%%%%%%%
\subsection{Cycle-Service Level}
\label{subsec:cycle-service-level}
%%%%%%%%%%%%%%%%%%%%%%%%%%%%%%%%%%%%%%%%%%%%%%%%%%%%%%%%%%%%

The cycle-service level is the probability that no stockout occurs during
a replenishment cycle.

In a basic continuous-review model,

\[
\boxed{
CSL
=
\mathbb{P}
(
D_L\leq r
).
}
\]

Thus choosing

\[
r
=
F_{D_L}^{-1}(\alpha)
\]

targets cycle-service level \(\alpha\).

%%%%%%%%%%%%%%%%%%%%%%%%%%%%%%%%%%%%%%%%%%%%%%%%%%%%%%%%%%%%
\subsection{Fill Rate}
\label{subsec:fill-rate}
%%%%%%%%%%%%%%%%%%%%%%%%%%%%%%%%%%%%%%%%%%%%%%%%%%%%%%%%%%%%

Fill rate measures the fraction of demand immediately satisfied from
available inventory.

A basic unit fill rate can be written

\[
\boxed{
\beta
=
1
-
\frac{
\text{expected shortage units per cycle}
}{
Q
}.
}
\]

Cycle-service level and fill rate are different performance measures.

%%%%%%%%%%%%%%%%%%%%%%%%%%%%%%%%%%%%%%%%%%%%%%%%%%%%%%%%%%%%
\subsection{Cycle Service Versus Fill Rate}
\label{subsec:cycle-service-versus-fill-rate}
%%%%%%%%%%%%%%%%%%%%%%%%%%%%%%%%%%%%%%%%%%%%%%%%%%%%%%%%%%%%

Cycle-service level asks:

\[
\boxed{
\text{Was there any stockout during the cycle?}
}
\]

Fill rate asks:

\[
\boxed{
\text{What fraction of total demand was immediately filled?}
}
\]

A system may have occasional stockouts but still achieve a very high fill
rate if shortages are small.

%%%%%%%%%%%%%%%%%%%%%%%%%%%%%%%%%%%%%%%%%%%%%%%%%%%%%%%%%%%%
\subsection{Continuous-Review \((Q,r)\) Policy}
\label{subsec:q-r-policy}
%%%%%%%%%%%%%%%%%%%%%%%%%%%%%%%%%%%%%%%%%%%%%%%%%%%%%%%%%%%%

A continuous-review \((Q,r)\) policy operates as follows:

\[
\boxed{
\text{when inventory position reaches }r,
\text{ order }Q\text{ units}.
}
\]

Thus:

\[
Q
=
\text{order quantity},
\]

and

\[
r
=
\text{reorder point}.
\]

%%%%%%%%%%%%%%%%%%%%%%%%%%%%%%%%%%%%%%%%%%%%%%%%%%%%%%%%%%%%
\subsection{Separation of \(Q\) and \(r\)}
\label{subsec:q-r-separation}
%%%%%%%%%%%%%%%%%%%%%%%%%%%%%%%%%%%%%%%%%%%%%%%%%%%%%%%%%%%%

In many classical approximations:

\[
\boxed{
Q
}
\]

primarily balances ordering and holding costs,

while

\[
\boxed{
r
}
\]

primarily controls shortage risk during lead time.

This allows EOQ and safety-stock calculations to be combined.

%%%%%%%%%%%%%%%%%%%%%%%%%%%%%%%%%%%%%%%%%%%%%%%%%%%%%%%%%%%%
\subsection{Periodic Review}
\label{subsec:periodic-review}
%%%%%%%%%%%%%%%%%%%%%%%%%%%%%%%%%%%%%%%%%%%%%%%%%%%%%%%%%%%%

Under periodic review, inventory is examined every

\[
R
\]

time units.

At each review, an order may be placed to raise inventory position to a
target level

\[
\boxed{
S.
}
\]

This is an order-up-to policy.

%%%%%%%%%%%%%%%%%%%%%%%%%%%%%%%%%%%%%%%%%%%%%%%%%%%%%%%%%%%%
\subsection{Order-Up-To Policy}
\label{subsec:order-up-to-policy}
%%%%%%%%%%%%%%%%%%%%%%%%%%%%%%%%%%%%%%%%%%%%%%%%%%%%%%%%%%%%

Suppose current inventory position is \(IP\).

Then the order quantity is

\[
\boxed{
Q
=
S-IP,
}
\]

if this value is positive.

The target \(S\) must protect against demand over both:

\[
\boxed{
\text{review period}
}
\]

and

\[
\boxed{
\text{lead time}.
}
\]

%%%%%%%%%%%%%%%%%%%%%%%%%%%%%%%%%%%%%%%%%%%%%%%%%%%%%%%%%%%%
\subsection{Protection Period}
\label{subsec:inventory-protection-period}
%%%%%%%%%%%%%%%%%%%%%%%%%%%%%%%%%%%%%%%%%%%%%%%%%%%%%%%%%%%%

For a periodic-review system with review period \(R\) and lead time \(L\),
the protection period is

\[
\boxed{
R+L.
}
\]

The order-up-to level should therefore account for demand over \(R+L\),
not merely during lead time.

%%%%%%%%%%%%%%%%%%%%%%%%%%%%%%%%%%%%%%%%%%%%%%%%%%%%%%%%%%%%
\subsection{Periodic-Review Normal Approximation}
\label{subsec:periodic-review-normal}
%%%%%%%%%%%%%%%%%%%%%%%%%%%%%%%%%%%%%%%%%%%%%%%%%%%%%%%%%%%%

If demand over the protection period is approximately normal with mean

\[
\mu_{R+L}
\]

and standard deviation

\[
\sigma_{R+L},
\]

a service-based order-up-to level is

\[
\boxed{
S
=
\mu_{R+L}
+
z_\alpha
\sigma_{R+L}.
}
\]

%%%%%%%%%%%%%%%%%%%%%%%%%%%%%%%%%%%%%%%%%%%%%%%%%%%%%%%%%%%%
\subsection{Continuous Versus Periodic Review}
\label{subsec:continuous-versus-periodic-inventory}
%%%%%%%%%%%%%%%%%%%%%%%%%%%%%%%%%%%%%%%%%%%%%%%%%%%%%%%%%%%%

Continuous review can react immediately when inventory position reaches a
threshold.

Periodic review waits until the next scheduled review.

Therefore periodic review usually needs protection against uncertainty
over a longer interval.

%%%%%%%%%%%%%%%%%%%%%%%%%%%%%%%%%%%%%%%%%%%%%%%%%%%%%%%%%%%%
\subsection{The Newsvendor Problem}
\label{subsec:newsvendor-problem}
%%%%%%%%%%%%%%%%%%%%%%%%%%%%%%%%%%%%%%%%%%%%%%%%%%%%%%%%%%%%

The newsvendor problem models a one-period inventory decision under
uncertain demand.

Choose order quantity

\[
\boxed{
Q
}
\]

before random demand

\[
\boxed{
D
}
\]

is observed.

Ordering too little causes shortage cost.

Ordering too much causes leftover inventory cost.

%%%%%%%%%%%%%%%%%%%%%%%%%%%%%%%%%%%%%%%%%%%%%%%%%%%%%%%%%%%%
\subsection{Overage and Underage Costs}
\label{subsec:overage-underage-costs}
%%%%%%%%%%%%%%%%%%%%%%%%%%%%%%%%%%%%%%%%%%%%%%%%%%%%%%%%%%%%

Define:

\[
\boxed{
C_u
=
\text{underage cost per unit},
}
\]

and

\[
\boxed{
C_o
=
\text{overage cost per unit}.
}
\]

Underage cost measures the marginal loss from ordering one unit too few.

Overage cost measures the marginal loss from ordering one unit too many.

%%%%%%%%%%%%%%%%%%%%%%%%%%%%%%%%%%%%%%%%%%%%%%%%%%%%%%%%%%%%
\subsection{Newsvendor Critical Fractile}
\label{subsec:newsvendor-critical-fractile}
%%%%%%%%%%%%%%%%%%%%%%%%%%%%%%%%%%%%%%%%%%%%%%%%%%%%%%%%%%%%

For continuous demand distribution function \(F_D\), the optimal order
quantity satisfies

\[
\boxed{
F_D(Q^*)
=
\frac{
C_u
}{
C_u+C_o
}.
}
\]

Therefore,

\[
\boxed{
Q^*
=
F_D^{-1}
\left(
\frac{
C_u
}{
C_u+C_o
}
\right).
}
\]

This is the critical-fractile solution.

%%%%%%%%%%%%%%%%%%%%%%%%%%%%%%%%%%%%%%%%%%%%%%%%%%%%%%%%%%%%
\subsection{Derivation of the Critical Fractile}
\label{subsec:newsvendor-critical-fractile-derivation}
%%%%%%%%%%%%%%%%%%%%%%%%%%%%%%%%%%%%%%%%%%%%%%%%%%%%%%%%%%%%

Consider increasing \(Q\) by one small unit.

The extra unit is useful when

\[
D>Q,
\]

which occurs with probability

\[
1-F_D(Q).
\]

Its expected marginal benefit is approximately

\[
C_u
[
1-F_D(Q)
].
\]

The unit is left over when

\[
D\leq Q,
\]

with probability

\[
F_D(Q).
\]

Its expected marginal cost is approximately

\[
C_oF_D(Q).
\]

At the optimum,

\[
C_u
[
1-F_D(Q)
]
=
C_oF_D(Q).
\]

Therefore,

\[
\boxed{
F_D(Q)
=
\frac{
C_u
}{
C_u+C_o
}.
}
\]

%%%%%%%%%%%%%%%%%%%%%%%%%%%%%%%%%%%%%%%%%%%%%%%%%%%%%%%%%%%%
\subsection{Newsvendor With Selling and Salvage Values}
\label{subsec:newsvendor-selling-salvage}
%%%%%%%%%%%%%%%%%%%%%%%%%%%%%%%%%%%%%%%%%%%%%%%%%%%%%%%%%%%%

Suppose:

\[
c
=
\text{purchase cost},
\]

\[
p
=
\text{selling price},
\]

and

\[
v
=
\text{salvage value},
\]

with

\[
p>c>v.
\]

Then a common formulation gives

\[
\boxed{
C_u
=
p-c,
}
\]

and

\[
\boxed{
C_o
=
c-v.
}
\]

Thus,

\[
\boxed{
F_D(Q^*)
=
\frac{
p-c
}{
p-v
}.
}
\]

%%%%%%%%%%%%%%%%%%%%%%%%%%%%%%%%%%%%%%%%%%%%%%%%%%%%%%%%%%%%
\subsection{Worked Example: Newsvendor}
\label{subsec:worked-newsvendor}
%%%%%%%%%%%%%%%%%%%%%%%%%%%%%%%%%%%%%%%%%%%%%%%%%%%%%%%%%%%%

\begin{workedexamplebox}{One-Period Order Quantity}

Suppose:

\[
p=20,
\qquad
c=12,
\qquad
v=4.
\]

Then

\[
C_u
=
20-12
=
8,
\]

and

\[
C_o
=
12-4
=
8.
\]

The critical fractile is

\[
\frac{
8
}{
8+8
}
=
0.5.
\]

Therefore,

\begin{answerbox}
\[
\boxed{
Q^*
=
F_D^{-1}(0.5),
}
\]

so the optimal order quantity is a median of demand under the standard
continuous formulation.
\end{answerbox}

\end{workedexamplebox}

%%%%%%%%%%%%%%%%%%%%%%%%%%%%%%%%%%%%%%%%%%%%%%%%%%%%%%%%%%%%
\subsection{Normal Newsvendor}
\label{subsec:normal-newsvendor}
%%%%%%%%%%%%%%%%%%%%%%%%%%%%%%%%%%%%%%%%%%%%%%%%%%%%%%%%%%%%

Suppose

\[
D
\sim
N(\mu,\sigma^2).
\]

Let

\[
\alpha
=
\frac{
C_u
}{
C_u+C_o
}.
\]

Then

\[
\boxed{
Q^*
=
\mu
+
z_\alpha\sigma.
}
\]

Thus the cost tradeoff determines the appropriate demand quantile.

%%%%%%%%%%%%%%%%%%%%%%%%%%%%%%%%%%%%%%%%%%%%%%%%%%%%%%%%%%%%
\subsection{Discrete Newsvendor}
\label{subsec:discrete-newsvendor}
%%%%%%%%%%%%%%%%%%%%%%%%%%%%%%%%%%%%%%%%%%%%%%%%%%%%%%%%%%%%

For discrete demand, the optimal quantity can be chosen as the smallest
integer \(Q\) satisfying

\[
\boxed{
F_D(Q)
\geq
\frac{
C_u
}{
C_u+C_o
}.
}
\]

Care is needed at probability jumps because the continuous equality need
not hold exactly.

%%%%%%%%%%%%%%%%%%%%%%%%%%%%%%%%%%%%%%%%%%%%%%%%%%%%%%%%%%%%
\subsection{Expected Shortage}
\label{subsec:expected-shortage}
%%%%%%%%%%%%%%%%%%%%%%%%%%%%%%%%%%%%%%%%%%%%%%%%%%%%%%%%%%%%

For order quantity \(Q\), shortage is

\[
\boxed{
(D-Q)_+
=
\max
\{
D-Q,0
\}.
}
\]

Expected shortage is

\[
\boxed{
\mathbb{E}
[
(D-Q)_+
].
}
\]

Expected leftover inventory is

\[
\boxed{
\mathbb{E}
[
(Q-D)_+
].
}
\]

These quantities are useful for service and cost analysis.

%%%%%%%%%%%%%%%%%%%%%%%%%%%%%%%%%%%%%%%%%%%%%%%%%%%%%%%%%%%%
\subsection{Backorders}
\label{subsec:inventory-backorders}
%%%%%%%%%%%%%%%%%%%%%%%%%%%%%%%%%%%%%%%%%%%%%%%%%%%%%%%%%%%%

If unmet demand is backordered, customers wait until inventory becomes
available.

A backorder quantity is therefore negative net inventory in a common
inventory-level convention.

Future replenishment first satisfies accumulated backorders before
building positive on-hand inventory.

%%%%%%%%%%%%%%%%%%%%%%%%%%%%%%%%%%%%%%%%%%%%%%%%%%%%%%%%%%%%
\subsection{Lost Sales}
\label{subsec:lost-sales}
%%%%%%%%%%%%%%%%%%%%%%%%%%%%%%%%%%%%%%%%%%%%%%%%%%%%%%%%%%%%

If unmet demand leaves permanently, the demand is lost rather than
backordered.

This changes both:

\[
\boxed{
\text{future inventory dynamics}
}
\]

and

\[
\boxed{
\text{shortage cost}.
}
\]

Lost-sales inventory models are often more difficult than corresponding
backorder models.

%%%%%%%%%%%%%%%%%%%%%%%%%%%%%%%%%%%%%%%%%%%%%%%%%%%%%%%%%%%%
\subsection{Base-Stock Policy}
\label{subsec:base-stock-policy}
%%%%%%%%%%%%%%%%%%%%%%%%%%%%%%%%%%%%%%%%%%%%%%%%%%%%%%%%%%%%

A base-stock policy attempts to restore inventory position to a fixed
target

\[
\boxed{
S.
}
\]

After demand of \(d\) units reduces inventory position by \(d\), the policy
orders approximately \(d\) units to restore the target.

Such policies arise naturally when fixed ordering costs are absent.

%%%%%%%%%%%%%%%%%%%%%%%%%%%%%%%%%%%%%%%%%%%%%%%%%%%%%%%%%%%%
\subsection{\((s,S)\) Policy}
\label{subsec:s-S-policy}
%%%%%%%%%%%%%%%%%%%%%%%%%%%%%%%%%%%%%%%%%%%%%%%%%%%%%%%%%%%%

An \((s,S)\) policy operates as follows:

\[
\boxed{
\text{if inventory position}\leq s,
\text{ order enough to raise it to }S.
}
\]

Thus order quantity is

\[
\boxed{
S-IP
}
\]

when an order is triggered.

Fixed ordering costs often create this threshold behavior.

%%%%%%%%%%%%%%%%%%%%%%%%%%%%%%%%%%%%%%%%%%%%%%%%%%%%%%%%%%%%
\subsection{Difference Between \((Q,r)\) and \((s,S)\)}
\label{subsec:q-r-versus-s-S}
%%%%%%%%%%%%%%%%%%%%%%%%%%%%%%%%%%%%%%%%%%%%%%%%%%%%%%%%%%%%

A \((Q,r)\) policy orders a fixed amount \(Q\):

\[
\boxed{
IP\leq r
\Rightarrow
\text{order }Q.
}
\]

An \((s,S)\) policy orders a variable quantity:

\[
\boxed{
IP\leq s
\Rightarrow
\text{raise }IP\text{ to }S.
}
\]

The two policies can behave differently under variable demand.

%%%%%%%%%%%%%%%%%%%%%%%%%%%%%%%%%%%%%%%%%%%%%%%%%%%%%%%%%%%%
\subsection{Finite-Horizon Inventory Dynamic Programming}
\label{subsec:finite-horizon-inventory-dp}
%%%%%%%%%%%%%%%%%%%%%%%%%%%%%%%%%%%%%%%%%%%%%%%%%%%%%%%%%%%%

Let:

\[
x_t
=
\text{inventory before ordering at stage }t,
\]

\[
q_t
=
\text{order quantity},
\]

and

\[
D_t
=
\text{random demand}.
\]

A simple transition is

\[
\boxed{
x_{t+1}
=
x_t+q_t-D_t.
}
\]

%%%%%%%%%%%%%%%%%%%%%%%%%%%%%%%%%%%%%%%%%%%%%%%%%%%%%%%%%%%%
\subsection{Inventory Bellman Equation}
\label{subsec:inventory-bellman-equation}
%%%%%%%%%%%%%%%%%%%%%%%%%%%%%%%%%%%%%%%%%%%%%%%%%%%%%%%%%%%%

Suppose stage cost is

\[
c_t(x_t,q_t,D_t).
\]

Then

\[
\boxed{
V_t(x)
=
\min_{q\geq0}
\mathbb{E}
\left[
c_t(x,q,D_t)
+
V_{t+1}
(
x+q-D_t
)
\right].
}
\]

This is a stochastic dynamic program.

%%%%%%%%%%%%%%%%%%%%%%%%%%%%%%%%%%%%%%%%%%%%%%%%%%%%%%%%%%%%
\subsection{Fixed Ordering Cost in Dynamic Inventory}
\label{subsec:fixed-ordering-cost-dynamic}
%%%%%%%%%%%%%%%%%%%%%%%%%%%%%%%%%%%%%%%%%%%%%%%%%%%%%%%%%%%%

Suppose placing any positive order costs

\[
K.
\]

Then an ordering cost may be written

\[
\boxed{
K\mathbf{1}_{\{q>0\}}
+
cq,
}
\]

where

\[
\mathbf{1}_{\{q>0\}}
=
\begin{cases}
1,
&
q>0,
\\
0,
&
q=0.
\end{cases}
\]

The fixed charge encourages batching and can lead to \((s,S)\)-type
policies.

%%%%%%%%%%%%%%%%%%%%%%%%%%%%%%%%%%%%%%%%%%%%%%%%%%%%%%%%%%%%
\subsection{Holding and Shortage Cost Function}
\label{subsec:inventory-holding-shortage-function}
%%%%%%%%%%%%%%%%%%%%%%%%%%%%%%%%%%%%%%%%%%%%%%%%%%%%%%%%%%%%

A common end-of-period cost is

\[
\boxed{
h(x)^+
+
p(-x)^+,
}
\]

where

\[
(x)^+
=
\max\{x,0\}.
\]

Thus positive ending inventory incurs holding cost, while negative
inventory incurs shortage or backorder cost.

%%%%%%%%%%%%%%%%%%%%%%%%%%%%%%%%%%%%%%%%%%%%%%%%%%%%%%%%%%%%
\subsection{Convexity in Inventory Models}
\label{subsec:inventory-convexity}
%%%%%%%%%%%%%%%%%%%%%%%%%%%%%%%%%%%%%%%%%%%%%%%%%%%%%%%%%%%%

Many inventory cost functions are convex.

For example,

\[
\boxed{
h(x)^+
+
p(-x)^+
}
\]

is convex in \(x\).

Convexity helps explain threshold and order-up-to policy structures.

%%%%%%%%%%%%%%%%%%%%%%%%%%%%%%%%%%%%%%%%%%%%%%%%%%%%%%%%%%%%
\subsection{Newsvendor as a One-Period Dynamic Program}
\label{subsec:newsvendor-one-period-dp}
%%%%%%%%%%%%%%%%%%%%%%%%%%%%%%%%%%%%%%%%%%%%%%%%%%%%%%%%%%%%

The newsvendor model can be viewed as the terminal-stage version of a
multi-period inventory problem.

There is:

\[
\boxed{
\text{one order decision},
}
\]

followed by

\[
\boxed{
\text{one uncertain demand realization}.
}
\]

This makes newsvendor theory foundational for stochastic inventory
control.

%%%%%%%%%%%%%%%%%%%%%%%%%%%%%%%%%%%%%%%%%%%%%%%%%%%%%%%%%%%%
\subsection{Multi-Item Inventory}
\label{subsec:multi-item-inventory}
%%%%%%%%%%%%%%%%%%%%%%%%%%%%%%%%%%%%%%%%%%%%%%%%%%%%%%%%%%%%

Organizations usually manage many items simultaneously.

If item \(i\) has order quantity \(Q_i\), a joint capacity constraint may
be

\[
\boxed{
\sum_i
a_iQ_i
\leq
B.
}
\]

Here \(a_i\) could represent storage space or capital required per unit.

%%%%%%%%%%%%%%%%%%%%%%%%%%%%%%%%%%%%%%%%%%%%%%%%%%%%%%%%%%%%
\subsection{Budget-Constrained Inventory}
\label{subsec:budget-constrained-inventory}
%%%%%%%%%%%%%%%%%%%%%%%%%%%%%%%%%%%%%%%%%%%%%%%%%%%%%%%%%%%%

Suppose unit value of item \(i\) is \(c_i\).

A budget constraint may be

\[
\boxed{
\sum_i
c_i
I_i
\leq
B,
}
\]

where \(I_i\) is an appropriate inventory measure.

The individual EOQ solutions may no longer be jointly feasible.

%%%%%%%%%%%%%%%%%%%%%%%%%%%%%%%%%%%%%%%%%%%%%%%%%%%%%%%%%%%%
\subsection{ABC Classification}
\label{subsec:abc-inventory}
%%%%%%%%%%%%%%%%%%%%%%%%%%%%%%%%%%%%%%%%%%%%%%%%%%%%%%%%%%%%

ABC classification prioritizes inventory-management effort.

A typical interpretation is:

\[
\boxed{
A
=
\text{few high-value or high-impact items},
}
\]

\[
\boxed{
B
=
\text{intermediate items},
}
\]

and

\[
\boxed{
C
=
\text{many lower-impact items}.
}
\]

The exact thresholds depend on the organization.

%%%%%%%%%%%%%%%%%%%%%%%%%%%%%%%%%%%%%%%%%%%%%%%%%%%%%%%%%%%%
\subsection{Annual Dollar Usage}
\label{subsec:annual-dollar-usage}
%%%%%%%%%%%%%%%%%%%%%%%%%%%%%%%%%%%%%%%%%%%%%%%%%%%%%%%%%%%%

A common classification measure is

\[
\boxed{
\text{annual dollar usage}
=
\text{annual demand}
\times
\text{unit value}.
}
\]

Items can be ranked by this measure before assigning ABC categories.

ABC analysis is a managerial prioritization tool rather than an
optimization theorem.

%%%%%%%%%%%%%%%%%%%%%%%%%%%%%%%%%%%%%%%%%%%%%%%%%%%%%%%%%%%%
\subsection{Perishable Inventory}
\label{subsec:perishable-inventory}
%%%%%%%%%%%%%%%%%%%%%%%%%%%%%%%%%%%%%%%%%%%%%%%%%%%%%%%%%%%%

Some inventory deteriorates or expires.

Examples include:

\[
\boxed{
\text{food},
}
\]

\[
\boxed{
\text{medicine},
}
\]

and

\[
\boxed{
\text{time-sensitive products}.
}
\]

Age or remaining shelf life then becomes part of the inventory state.

%%%%%%%%%%%%%%%%%%%%%%%%%%%%%%%%%%%%%%%%%%%%%%%%%%%%%%%%%%%%
\subsection{Spare-Parts Inventory}
\label{subsec:spare-parts-inventory}
%%%%%%%%%%%%%%%%%%%%%%%%%%%%%%%%%%%%%%%%%%%%%%%%%%%%%%%%%%%%

Spare-parts systems may have:

\[
\boxed{
\text{low demand rates},
}
\]

\[
\boxed{
\text{high shortage costs},
}
\]

and

\[
\boxed{
\text{long replenishment lead times}.
}
\]

Demand may be intermittent rather than well approximated by ordinary
continuous distributions.

%%%%%%%%%%%%%%%%%%%%%%%%%%%%%%%%%%%%%%%%%%%%%%%%%%%%%%%%%%%%
\subsection{Service Parts and Availability}
\label{subsec:service-parts-availability}
%%%%%%%%%%%%%%%%%%%%%%%%%%%%%%%%%%%%%%%%%%%%%%%%%%%%%%%%%%%%

For maintenance systems, inventory performance may be measured through
equipment availability rather than ordinary sales fill rate.

A missing critical spare part can keep expensive equipment out of service.

Thus shortage costs may be highly asymmetric across items.

%%%%%%%%%%%%%%%%%%%%%%%%%%%%%%%%%%%%%%%%%%%%%%%%%%%%%%%%%%%%
\subsection{Multi-Echelon Inventory}
\label{subsec:multi-echelon-inventory}
%%%%%%%%%%%%%%%%%%%%%%%%%%%%%%%%%%%%%%%%%%%%%%%%%%%%%%%%%%%%

A multi-echelon inventory system contains several levels such as

\[
\boxed{
\text{supplier}
\to
\text{distribution center}
\to
\text{retailer}.
}
\]

Inventory decisions at one level affect replenishment and service at other
levels.

Local optimization of each echelon may not minimize total system cost.

%%%%%%%%%%%%%%%%%%%%%%%%%%%%%%%%%%%%%%%%%%%%%%%%%%%%%%%%%%%%
\subsection{Echelon Inventory Position}
\label{subsec:echelon-inventory-position}
%%%%%%%%%%%%%%%%%%%%%%%%%%%%%%%%%%%%%%%%%%%%%%%%%%%%%%%%%%%%

Echelon inventory includes inventory at a location and downstream
locations, adjusted for downstream demand commitments according to the
chosen convention.

This broader state variable helps coordinate multi-stage supply chains.

%%%%%%%%%%%%%%%%%%%%%%%%%%%%%%%%%%%%%%%%%%%%%%%%%%%%%%%%%%%%
\subsection{Risk Pooling}
\label{subsec:inventory-risk-pooling}
%%%%%%%%%%%%%%%%%%%%%%%%%%%%%%%%%%%%%%%%%%%%%%%%%%%%%%%%%%%%

Combining demand from several locations can reduce relative uncertainty.

If independent demands have standard deviations

\[
\sigma_1,\ldots,\sigma_n,
\]

then the standard deviation of aggregate demand is

\[
\boxed{
\sqrt{
\sum_{i=1}^{n}
\sigma_i^2
}.
}
\]

This is generally less than

\[
\sum_i\sigma_i.
\]

Pooling can therefore reduce total safety-stock requirements under
appropriate conditions.

%%%%%%%%%%%%%%%%%%%%%%%%%%%%%%%%%%%%%%%%%%%%%%%%%%%%%%%%%%%%
\subsection{Correlation and Risk Pooling}
\label{subsec:correlation-risk-pooling}
%%%%%%%%%%%%%%%%%%%%%%%%%%%%%%%%%%%%%%%%%%%%%%%%%%%%%%%%%%%%

If demands are correlated,

\[
\operatorname{Var}
\left(
\sum_iD_i
\right)
=
\sum_i
\operatorname{Var}(D_i)
+
2
\sum_{i<j}
\operatorname{Cov}(D_i,D_j).
\]

Positive correlation reduces the benefit of pooling.

Negative correlation can increase it.

Thus diversification benefits depend on dependence structure.

%%%%%%%%%%%%%%%%%%%%%%%%%%%%%%%%%%%%%%%%%%%%%%%%%%%%%%%%%%%%
\subsection{Bullwhip Effect}
\label{subsec:bullwhip-effect}
%%%%%%%%%%%%%%%%%%%%%%%%%%%%%%%%%%%%%%%%%%%%%%%%%%%%%%%%%%%%

The bullwhip effect occurs when demand variability becomes amplified as
one moves upstream through a supply chain.

Small fluctuations in customer demand may produce larger fluctuations in:

\[
\boxed{
\text{retailer orders},
}
\]

\[
\boxed{
\text{wholesaler orders},
}
\]

and

\[
\boxed{
\text{manufacturer production}.
}
\]

%%%%%%%%%%%%%%%%%%%%%%%%%%%%%%%%%%%%%%%%%%%%%%%%%%%%%%%%%%%%
\subsection{Causes of the Bullwhip Effect}
\label{subsec:bullwhip-causes}
%%%%%%%%%%%%%%%%%%%%%%%%%%%%%%%%%%%%%%%%%%%%%%%%%%%%%%%%%%%%

Potential causes include:

\[
\boxed{
\text{forecast updating},
}
\]

\[
\boxed{
\text{batch ordering},
}
\]

\[
\boxed{
\text{price promotions},
}
\]

and

\[
\boxed{
\text{shortage gaming}.
}
\]

Information sharing and coordinated replenishment can reduce amplification.

%%%%%%%%%%%%%%%%%%%%%%%%%%%%%%%%%%%%%%%%%%%%%%%%%%%%%%%%%%%%
\subsection{Inventory and Queueing Connection}
\label{subsec:inventory-queueing-connection}
%%%%%%%%%%%%%%%%%%%%%%%%%%%%%%%%%%%%%%%%%%%%%%%%%%%%%%%%%%%%

Inventory and queueing systems are dual in an operational sense.

In inventory:

\[
\boxed{
\text{supply waits for demand}.
}
\]

In queueing:

\[
\boxed{
\text{demand waits for supply or service}.
}
\]

Both depend strongly on:

\[
\boxed{
\text{rates},
}
\]

\[
\boxed{
\text{variability},
}
\]

and

\[
\boxed{
\text{capacity}.
}
\]

%%%%%%%%%%%%%%%%%%%%%%%%%%%%%%%%%%%%%%%%%%%%%%%%%%%%%%%%%%%%
\subsection{Little's Law and Pipeline Inventory}
\label{subsec:littles-law-pipeline-inventory}
%%%%%%%%%%%%%%%%%%%%%%%%%%%%%%%%%%%%%%%%%%%%%%%%%%%%%%%%%%%%

Little's Law can also interpret material moving through a supply system.

If throughput is

\[
\lambda
\]

and average replenishment lead time is

\[
W,
\]

then average pipeline inventory is approximately

\[
\boxed{
L
=
\lambda W.
}
\]

Thus longer lead times create more inventory tied up in transit.

%%%%%%%%%%%%%%%%%%%%%%%%%%%%%%%%%%%%%%%%%%%%%%%%%%%%%%%%%%%%
\subsection{Inventory and Scheduling}
\label{subsec:inventory-scheduling-connection}
%%%%%%%%%%%%%%%%%%%%%%%%%%%%%%%%%%%%%%%%%%%%%%%%%%%%%%%%%%%%

Production scheduling determines when items are produced.

Inventory determines how much material should be available before and
after production.

These decisions interact:

\[
\boxed{
\text{larger production batches}
\rightarrow
\text{more cycle inventory},
}
\]

while

\[
\boxed{
\text{smaller batches}
\rightarrow
\text{more frequent setups}.
}
\]

This is the production analogue of EOQ tradeoffs.

%%%%%%%%%%%%%%%%%%%%%%%%%%%%%%%%%%%%%%%%%%%%%%%%%%%%%%%%%%%%
\subsection{Inventory and Stochastic Optimization}
\label{subsec:inventory-stochastic-optimization}
%%%%%%%%%%%%%%%%%%%%%%%%%%%%%%%%%%%%%%%%%%%%%%%%%%%%%%%%%%%%

When demand is uncertain, inventory decisions solve stochastic
optimization problems.

Examples include:

\[
\boxed{
\text{newsvendor},
}
\]

\[
\boxed{
(Q,r)\text{ optimization},
}
\]

and

\[
\boxed{
\text{multiperiod inventory control}.
}
\]

These models use probability distributions, expectations, and service
constraints.

%%%%%%%%%%%%%%%%%%%%%%%%%%%%%%%%%%%%%%%%%%%%%%%%%%%%%%%%%%%%
\subsection{Inventory Optimization Workflow}
\label{subsec:inventory-workflow}
%%%%%%%%%%%%%%%%%%%%%%%%%%%%%%%%%%%%%%%%%%%%%%%%%%%%%%%%%%%%

A practical workflow is:

\begin{enumerate}
    \item identify the inventory item;
    \item estimate demand rate or demand distribution;
    \item determine replenishment lead time;
    \item determine whether review is continuous or periodic;
    \item estimate ordering or setup cost;
    \item estimate holding cost;
    \item identify shortage behavior;
    \item determine whether shortages are backordered or lost;
    \item define the desired service measure;
    \item choose a suitable inventory model;
    \item compute economic order or production quantity when appropriate;
    \item compute reorder points or order-up-to levels;
    \item calculate safety stock;
    \item test sensitivity to demand and lead-time assumptions;
    \item verify units and time scales;
    \item consider multi-item or capacity constraints;
    \item evaluate total expected cost and service;
    \item update parameters as new data become available.
\end{enumerate}

%%%%%%%%%%%%%%%%%%%%%%%%%%%%%%%%%%%%%%%%%%%%%%%%%%%%%%%%%%%%
\subsection{Choosing an Inventory Model}
\label{subsec:choosing-inventory-model}
%%%%%%%%%%%%%%%%%%%%%%%%%%%%%%%%%%%%%%%%%%%%%%%%%%%%%%%%%%%%

For deterministic constant demand with instantaneous replenishment, use:

\[
\boxed{
\text{EOQ}.
}
\]

For finite production rate, use:

\[
\boxed{
\text{EPQ}.
}
\]

For one uncertain selling period, use:

\[
\boxed{
\text{newsvendor}.
}
\]

For continuous review with uncertain lead-time demand, consider:

\[
\boxed{
(Q,r).
}
\]

For periodic review, consider:

\[
\boxed{
\text{order-up-to }S.
}
\]

For multiperiod stochastic systems, use:

\[
\boxed{
\text{dynamic programming}
}
\]

or other stochastic optimization methods.

%%%%%%%%%%%%%%%%%%%%%%%%%%%%%%%%%%%%%%%%%%%%%%%%%%%%%%%%%%%%
\subsection{Common Errors}
\label{subsec:inventory-common-errors}
%%%%%%%%%%%%%%%%%%%%%%%%%%%%%%%%%%%%%%%%%%%%%%%%%%%%%%%%%%%%

\subsubsection{Confusing Demand Rate and Total Demand}

If

\[
D
\]

is annual demand, units must remain consistent with annual holding and
ordering costs.

Mixing daily and annual units produces incorrect results.

%%%%%%%%%%%%%%%%%%%%%%%%%%%%%%%%%%%%%%%%%%%%%%%%%%%%%%%%%%%%
\subsubsection{Forgetting the Square Root in EOQ}

The EOQ formula is

\[
\boxed{
Q^*
=
\sqrt{
\frac{2KD}{h}
}.
}
\]

It is not

\[
\frac{2KD}{h}.
\]

%%%%%%%%%%%%%%%%%%%%%%%%%%%%%%%%%%%%%%%%%%%%%%%%%%%%%%%%%%%%
\subsubsection{Including Constant Purchase Cost in EOQ Differentiation}

If unit purchase price does not depend on \(Q\), then

\[
cD
\]

is constant with respect to \(Q\).

It does not affect the EOQ minimizer.

It must be included when comparing price levels under quantity discounts.

%%%%%%%%%%%%%%%%%%%%%%%%%%%%%%%%%%%%%%%%%%%%%%%%%%%%%%%%%%%%
\subsubsection{Using \(Q/2\) With Finite Production Without Adjustment}

For EPQ,

\[
\boxed{
I_{\mathrm{avg}}
=
\frac{Q}{2}
\left(
1-\frac{d}{p}
\right),
}
\]

not simply \(Q/2\).

%%%%%%%%%%%%%%%%%%%%%%%%%%%%%%%%%%%%%%%%%%%%%%%%%%%%%%%%%%%%
\subsubsection{Using EOQ When Replenishment Is Not Instantaneous}

When production and demand occur simultaneously, EPQ may be more
appropriate.

The assumptions should match the physical system.

%%%%%%%%%%%%%%%%%%%%%%%%%%%%%%%%%%%%%%%%%%%%%%%%%%%%%%%%%%%%
\subsubsection{Confusing Reorder Point With Order Quantity}

The reorder point determines:

\[
\boxed{
\text{when to order}.
}
\]

The order quantity determines:

\[
\boxed{
\text{how much to order}.
}
\]

They are different decisions.

%%%%%%%%%%%%%%%%%%%%%%%%%%%%%%%%%%%%%%%%%%%%%%%%%%%%%%%%%%%%
\subsubsection{Ignoring Lead Time}

If replenishment is delayed, ordering at zero inventory generally causes a
stockout.

The reorder point must account for lead-time demand.

%%%%%%%%%%%%%%%%%%%%%%%%%%%%%%%%%%%%%%%%%%%%%%%%%%%%%%%%%%%%
\subsubsection{Using Safety Stock Without a Service Interpretation}

Safety stock should be connected to a probability or cost objective.

Choosing an arbitrary number of extra units gives no explicit service
meaning.

%%%%%%%%%%%%%%%%%%%%%%%%%%%%%%%%%%%%%%%%%%%%%%%%%%%%%%%%%%%%
\subsubsection{Confusing Cycle-Service Level and Fill Rate}

A \(95\%\) cycle-service level does not mean exactly \(95\%\) of units are
filled immediately.

These are different service metrics.

%%%%%%%%%%%%%%%%%%%%%%%%%%%%%%%%%%%%%%%%%%%%%%%%%%%%%%%%%%%%
\subsubsection{Using \(\sqrt{L}\sigma\) Without Appropriate Independence Assumptions}

The formula

\[
\sigma_L
=
\sqrt{L}\sigma_D
\]

assumes independent equal-period demand with fixed lead time.

Correlation or random lead time changes the variance.

%%%%%%%%%%%%%%%%%%%%%%%%%%%%%%%%%%%%%%%%%%%%%%%%%%%%%%%%%%%%
\subsubsection{Using the Continuous Newsvendor Equality for Discrete Demand}

For discrete demand, the CDF jumps.

Choose an appropriate quantile using an inequality rather than expecting
exact equality.

%%%%%%%%%%%%%%%%%%%%%%%%%%%%%%%%%%%%%%%%%%%%%%%%%%%%%%%%%%%%
\subsubsection{Using the Wrong Newsvendor Costs}

The critical fractile depends on marginal underage and overage costs.

These are not necessarily identical to selling price and purchase cost
themselves.

%%%%%%%%%%%%%%%%%%%%%%%%%%%%%%%%%%%%%%%%%%%%%%%%%%%%%%%%%%%%
\subsubsection{Ignoring Outstanding Orders}

Replenishment decisions should often use inventory position rather than
on-hand inventory.

Otherwise duplicate orders can occur.

%%%%%%%%%%%%%%%%%%%%%%%%%%%%%%%%%%%%%%%%%%%%%%%%%%%%%%%%%%%%
\subsubsection{Treating Lost Sales as Backorders}

Lost demand disappears.

Backordered demand remains and must be filled later.

The inventory dynamics are different.

%%%%%%%%%%%%%%%%%%%%%%%%%%%%%%%%%%%%%%%%%%%%%%%%%%%%%%%%%%%%
\subsubsection{Optimizing Each Echelon Independently}

A locally optimal policy at each warehouse or retailer may be globally
poor.

Multi-echelon inventory requires coordination.

%%%%%%%%%%%%%%%%%%%%%%%%%%%%%%%%%%%%%%%%%%%%%%%%%%%%%%%%%%%%
\subsubsection{Ignoring Demand Correlation in Risk Pooling}

Pooling benefits depend on covariance.

Highly positively correlated demands provide less diversification.

%%%%%%%%%%%%%%%%%%%%%%%%%%%%%%%%%%%%%%%%%%%%%%%%%%%%%%%%%%%%
\subsubsection{Assuming More Inventory Is Always Safer}

More inventory reduces shortage risk but increases:

\[
\boxed{
\text{holding cost},
}
\]

\[
\boxed{
\text{obsolescence risk},
}
\]

and

\[
\boxed{
\text{capital tied up}.
}
\]

Inventory management is an optimization tradeoff.

%%%%%%%%%%%%%%%%%%%%%%%%%%%%%%%%%%%%%%%%%%%%%%%%%%%%%%%%%%%%
\subsection{Applications}
\label{subsec:inventory-applications}
%%%%%%%%%%%%%%%%%%%%%%%%%%%%%%%%%%%%%%%%%%%%%%%%%%%%%%%%%%%%

\begin{applicationbox}

\textbf{Retail.}
Stores choose replenishment quantities and safety stock to maintain
availability while limiting carrying cost.

\medskip

\textbf{Manufacturing.}
Raw materials, components, work-in-process, and finished products must be
coordinated with production schedules.

\medskip

\textbf{Aviation and Maintenance.}
Spare parts must be available for maintenance while many parts have high
purchase costs and intermittent demand.

\medskip

\textbf{Healthcare.}
Hospitals manage medication, blood products, medical supplies, and
perishable inventory under uncertain demand.

\medskip

\textbf{E-Commerce.}
Distribution centers balance service speed, stock availability, and
inventory investment.

\medskip

\textbf{Supply Chains.}
Multi-echelon inventory decisions coordinate suppliers, warehouses,
distribution centers, and customers.

\medskip

\textbf{Energy.}
Fuel, battery storage, and reserve capacity can have inventory-like
decision structures.

\medskip

\textbf{Emergency Management.}
Relief supplies and critical equipment must be positioned before uncertain
future demand occurs.

\end{applicationbox}

%%%%%%%%%%%%%%%%%%%%%%%%%%%%%%%%%%%%%%%%%%%%%%%%%%%%%%%%%%%%
\subsection{Exercises}
\label{subsec:inventory-exercises}
%%%%%%%%%%%%%%%%%%%%%%%%%%%%%%%%%%%%%%%%%%%%%%%%%%%%%%%%%%%%

\begin{exercise}[1]
Define cycle stock.
\end{exercise}

\begin{exercise}[1]
Define safety stock.
\end{exercise}

\begin{exercise}[1]
Define pipeline inventory.
\end{exercise}

\begin{exercise}[1]
Define ordering cost.
\end{exercise}

\begin{exercise}[1]
Define holding cost.
\end{exercise}

\begin{exercise}[1]
Define shortage cost.
\end{exercise}

\begin{exercise}[1]
Define lead time.
\end{exercise}

\begin{exercise}[1]
Define inventory position.
\end{exercise}

\begin{exercise}[1]
Define cycle-service level.
\end{exercise}

\begin{exercise}[1]
Define fill rate.
\end{exercise}

\begin{exercise}[2]
If

\[
D=1000,
\qquad
Q=100,
\]

compute the number of orders per year.
\end{exercise}

\begin{exercise}[2]
For \(Q=200\), compute average cycle inventory under the basic EOQ model.
\end{exercise}

\begin{exercise}[2]
Write the EOQ relevant-cost function.
\end{exercise}

\begin{exercise}[2]
Compute EOQ for

\[
D=5000,
\quad
K=40,
\quad
h=8.
\]
\end{exercise}

\begin{exercise}[2]
Compute order frequency and cycle length from the previous EOQ.
\end{exercise}

\begin{exercise}[2]
Verify that ordering and holding costs are equal at EOQ.
\end{exercise}

\begin{exercise}[3]
Derive the EOQ formula using first-order calculus.
\end{exercise}

\begin{exercise}[3]
Use the second derivative to verify that EOQ minimizes relevant cost.
\end{exercise}

\begin{exercise}[3]
Derive

\[
TC(Q^*)
=
\sqrt{2KDh}.
\]
\end{exercise}

\begin{exercise}[3]
Compare total relevant cost at EOQ with costs at quantities \(20\%\)
above and below EOQ.
\end{exercise}

\begin{exercise}[3]
Derive maximum inventory in the EPQ model.
\end{exercise}

\begin{exercise}[3]
Derive the EPQ formula.
\end{exercise}

\begin{exercise}[3]
Show that EOQ is the limit of EPQ as production rate approaches infinity.
\end{exercise}

\begin{exercise}[3]
Solve an EPQ problem with specified production and demand rates.
\end{exercise}

\begin{exercise}[3]
Solve an all-units quantity-discount problem with two price levels.
\end{exercise}

\begin{exercise}[3]
Explain why purchase cost must be included when comparing quantity
discount candidates.
\end{exercise}

\begin{exercise}[2]
If demand is \(30\) units per day and lead time is \(5\) days, compute the
deterministic reorder point.
\end{exercise}

\begin{exercise}[3]
Suppose daily demand has mean \(100\), standard deviation \(15\), and fixed
lead time \(4\) days. Compute lead-time mean and standard deviation.
\end{exercise}

\begin{exercise}[3]
Using the previous values and \(z=1.645\), compute safety stock and reorder
point.
\end{exercise}

\begin{exercise}[3]
Explain what the selected \(z\)-value means in terms of cycle-service
probability under the normal model.
\end{exercise}

\begin{exercise}[3]
Compare cycle-service level and fill rate.
\end{exercise}

\begin{exercise}[3]
Given expected shortage per cycle and \(Q\), compute the fill rate.
\end{exercise}

\begin{exercise}[3]
Formulate a \((Q,r)\) policy.
\end{exercise}

\begin{exercise}[3]
Formulate a periodic-review order-up-to policy.
\end{exercise}

\begin{exercise}[3]
Explain why a periodic-review system protects over \(R+L\).
\end{exercise}

\begin{exercise}[3]
Compute an order-up-to level under normally distributed protection-period
demand.
\end{exercise}

\begin{exercise}[2]
Define underage and overage cost in a newsvendor problem.
\end{exercise}

\begin{exercise}[3]
Derive the newsvendor critical fractile.
\end{exercise}

\begin{exercise}[3]
Suppose

\[
p=30,
\qquad
c=18,
\qquad
v=6.
\]

Compute \(C_u\), \(C_o\), and the critical fractile.
\end{exercise}

\begin{exercise}[3]
If demand is normally distributed with specified \(\mu\) and \(\sigma\),
compute the newsvendor order quantity.
\end{exercise}

\begin{exercise}[3]
Solve a discrete newsvendor problem from a probability table.
\end{exercise}

\begin{exercise}[3]
Compute expected shortage

\[
\mathbb{E}[(D-Q)_+]
\]

for a small discrete demand distribution.
\end{exercise}

\begin{exercise}[3]
Compute expected leftover inventory

\[
\mathbb{E}[(Q-D)_+].
\]
\end{exercise}

\begin{exercise}[3]
Explain the difference between lost sales and backorders.
\end{exercise}

\begin{exercise}[3]
Explain how fixed ordering costs motivate \((s,S)\)-type policies.
\end{exercise}

\begin{exercise}[3]
Write a finite-horizon inventory Bellman equation.
\end{exercise}

\begin{exercise}[3]
Model holding and shortage cost using positive-part functions.
\end{exercise}

\begin{exercise}[3]
Construct a two-period stochastic inventory problem.
\end{exercise}

\begin{exercise}[3]
Construct a multi-item storage-capacity constraint.
\end{exercise}

\begin{exercise}[3]
Rank several items by annual dollar usage for ABC analysis.
\end{exercise}

\begin{exercise}[3]
Explain why ABC classifications are managerial rather than universal
mathematical thresholds.
\end{exercise}

\begin{exercise}[3]
Explain risk pooling using two independent demand locations.
\end{exercise}

\begin{exercise}[3]
Compute aggregate demand variance when two demands are correlated.
\end{exercise}

\begin{exercise}[3]
Explain how positive correlation affects safety-stock pooling benefits.
\end{exercise}

\begin{exercise}[3]
Describe four causes of the bullwhip effect.
\end{exercise}

\begin{exercise}[3]
Use Little's Law to estimate pipeline inventory from throughput and lead
time.
\end{exercise}

\begin{exercise}[4]
Derive EOQ with planned backorders.
\end{exercise}

\begin{exercise}[4]
Study finite-production inventory with planned shortages.
\end{exercise}

\begin{exercise}[4]
Study incremental quantity discounts.
\end{exercise}

\begin{exercise}[4]
Derive reorder-point formulas when lead time is random.
\end{exercise}

\begin{exercise}[4]
Derive lead-time demand variance when both demand and lead time are random.
\end{exercise}

\begin{exercise}[4]
Study fill-rate optimization for continuous-review inventory.
\end{exercise}

\begin{exercise}[4]
Study the normal-loss function used in expected-shortage calculations.
\end{exercise}

\begin{exercise}[4]
Derive newsvendor optimality from expected profit differentiation.
\end{exercise}

\begin{exercise}[4]
Study risk-averse newsvendor models using CVaR.
\end{exercise}

\begin{exercise}[4]
Study multiperiod base-stock optimality under standard assumptions.
\end{exercise}

\begin{exercise}[4]
Study \((s,S)\) policy optimality with fixed ordering costs.
\end{exercise}

\begin{exercise}[4]
Investigate K-convexity in inventory theory.
\end{exercise}

\begin{exercise}[4]
Study lost-sales inventory dynamic programming.
\end{exercise}

\begin{exercise}[4]
Study perishable inventory with age-dependent states.
\end{exercise}

\begin{exercise}[4]
Study intermittent-demand spare-parts inventory.
\end{exercise}

\begin{exercise}[4]
Study multi-echelon base-stock systems.
\end{exercise}

\begin{exercise}[4]
Investigate Clark--Scarf decomposition for serial inventory systems.
\end{exercise}

\begin{exercise}[4]
Study risk pooling with correlated demands.
\end{exercise}

\begin{exercise}[4]
Develop a mathematical model of the bullwhip effect.
\end{exercise}

\begin{exercise}[4]
Investigate joint replenishment problems.
\end{exercise}

\begin{exercise}[4]
Study inventory-routing problems.
\end{exercise}

\begin{exercise}[4]
Study robust inventory optimization under uncertain demand distributions.
\end{exercise}

%%%%%%%%%%%%%%%%%%%%%%%%%%%%%%%%%%%%%%%%%%%%%%%%%%%%%%%%%%%%
\subsection{Conceptual Summary}
\label{subsec:inventory-summary}
%%%%%%%%%%%%%%%%%%%%%%%%%%%%%%%%%%%%%%%%%%%%%%%%%%%%%%%%%%%%

Inventory theory determines

\[
\boxed{
\text{how much to order}
}
\]

and

\[
\boxed{
\text{when to order}.
}
\]

The classical EOQ cost is

\[
\boxed{
TC(Q)
=
K\frac{D}{Q}
+
h\frac{Q}{2}.
}
\]

Its minimizer is

\[
\boxed{
Q^*
=
\sqrt{
\frac{2KD}{h}
}.
}
\]

With finite production rate,

\[
\boxed{
Q^*
=
\sqrt{
\frac{
2KD
}{
h
\left(
1-\frac{d}{p}
\right)
}
}.
}
\]

A deterministic reorder point is

\[
\boxed{
r=dL.
}
\]

With uncertain lead-time demand,

\[
\boxed{
r
=
\mathbb{E}[D_L]
+
SS.
}
\]

Under a normal model,

\[
\boxed{
SS
=
z_\alpha\sigma_L.
}
\]

The newsvendor critical fractile is

\[
\boxed{
F_D(Q^*)
=
\frac{
C_u
}{
C_u+C_o
}.
}
\]

Important replenishment policies include:

\[
\boxed{
(Q,r),
}
\]

\[
\boxed{
(s,S),
}
\]

and

\[
\boxed{
\text{periodic order-up-to policies}.
}
\]

Inventory theory therefore combines

\[
\boxed{
\text{optimization}
+
\text{probability}
+
\text{dynamic programming}
+
\text{service management}
+
\text{supply-chain planning}.
}
\]

%%%%%%%%%%%%%%%%%%%%%%%%%%%%%%%%%%%%%%%%%%%%%%%%%%%%%%%%%%%%
\subsection{Section Connections}
\label{subsec:inventory-connections}
%%%%%%%%%%%%%%%%%%%%%%%%%%%%%%%%%%%%%%%%%%%%%%%%%%%%%%%%%%%%

\chapterconnection{
Inventory Theory builds on the optimization framework of Chapter~10 and
the stochastic methods of Section~\ref{sec:stochastic-optimization}.
Classical EOQ and EPQ models balance setup and holding costs, while
reorder-point, safety-stock, and newsvendor models introduce uncertainty
and service requirements. Multiperiod inventory control connects directly
with Dynamic Programming, and pipeline inventory connects with Little's
Law from Queueing Theory. Inventory decisions also interact strongly with
Scheduling Theory, network logistics, and supply-chain coordination. The
next section, Decision Theory, develops mathematical frameworks for
choosing among alternatives under certainty, risk, uncertainty, multiple
criteria, and imperfect information.
}

%%%%%%%%%%%%%%%%%%%%%%%%%%%%%%%%%%%%%%%%%%%%%%%%%%%%%%%%%%%%
% SECTION SOURCE TRACKING
%%%%%%%%%%%%%%%%%%%%%%%%%%%%%%%%%%%%%%%%%%%%%%%%%%%%%%%%%%%%

%------------------------------------------------------------
% 10.13 Inventory Theory
%------------------------------------------------------------
%
% Source basis:
%   - Standard deterministic and stochastic inventory theory.
%   - Standard operations-research replenishment models.
%   - Standard dynamic inventory-control theory.
%
% Used for:
%   - inventory definitions and purposes;
%   - cycle, safety, pipeline, anticipation, and decoupling stock;
%   - ordering, holding, shortage, and purchase costs;
%   - deterministic and stochastic inventory;
%   - continuous and periodic review;
%   - EOQ;
%   - EOQ derivation and cost balance;
%   - cycle lengths and order frequency;
%   - EPQ;
%   - finite replenishment rates;
%   - quantity discounts;
%   - reorder points;
%   - inventory position;
%   - lead-time demand;
%   - safety stock;
%   - normal lead-time demand models;
%   - cycle-service levels;
%   - fill rates;
%   - continuous-review (Q,r) policies;
%   - periodic-review order-up-to policies;
%   - protection periods;
%   - newsvendor models;
%   - critical fractiles;
%   - overage and underage costs;
%   - continuous and discrete newsvendor models;
%   - expected shortages and leftovers;
%   - backorders and lost sales;
%   - base-stock and (s,S) policies;
%   - finite-horizon dynamic programming;
%   - convex inventory costs;
%   - multi-item constraints;
%   - ABC classification;
%   - perishable and spare-parts inventory;
%   - multi-echelon inventory;
%   - risk pooling;
%   - demand correlation;
%   - bullwhip effect;
%   - queueing and scheduling connections;
%   - applications and exercises.
%
% Notes:
%   - Decision Theory is developed next in Section 10.14.
%   - Game Theory follows in Section 10.15.
%   - Optimal Control follows in Section 10.16.
%   - Stochastic foundations are developed in Chapters 7 and 8.
%
%%%%%%%%%%%%%%%%%%%%%%%%%%%%%%%%%%%%%%%%%%%%%%%%%%%%%%%%%%%%